\begin{frame}{그림으로: 부분관측이 마르코프 성질을 깨는 방식}
  \begin{center}
  \includegraphics[width=0.72\linewidth]{ch03_hidden_state_speed.png}
  \end{center}
  \vskip0.2em
  \small
  같은 카메라 관측 $o_t$ 아래 ``빠르게 이동 중''과 ``느리게 이동 중''이라는
  서로 다른 실제 상태가 숨어 있으면, 같은 $o_t$ 에서 출발해도 다음
  프레임의 분포가 달라진다 --- 이것이 마르코프 성질이 깨지는 구체적
  방식.
\end{frame}
